%========================================================================================
\section{Intégrations externes}
\label{sec:integrations}
%========================================================================================

TravelHub utilise deux API publiques qui ne demandent pas de clé. Open-Meteo~\cite{openmeteo,openmeteogeo} donne le géocodage et la météo avec une interface REST. Countries~\cite{countries} décrit les pays avec une interface GraphQL. Tout le code d'intégration est dans \code{apps/integrations/services/}. Les clients renvoient des \code{dataclasses} immuables et jamais le JSON brut du fournisseur.

\subsection{Open-Meteo (REST)}

Pour avoir la météo d'une destination, il faut deux appels.

Le premier est le géocodage. Si une destination a été saisie sans coordonnées, on la cherche d'abord avec \code{geocoding-api.open-meteo.com/v1/search}, à partir de la ville et du code pays. Les coordonnées trouvées sont enregistrées dans la destination, ce qui évite de refaire cet appel ensuite.

Le deuxième appel, \code{api.open-meteo.com/v1/forecast}, renvoie les conditions actuelles et les prévisions des 16 prochains jours. Le service garde seulement les jours du séjour.

Les codes météo WMO sont traduits en français (\og{}Pluie modérée\fg{}, \og{}Orage\fg{}\ldots). Si le séjour est complètement en dehors des 16 jours, on renvoie l'erreur \code{no_data_available}. Nous gardons un jour de marge, car la date du serveur et la date locale de la destination ne sont pas toujours les mêmes.

Pour enregistrer les coordonnées, nous utilisons \code{update()} avec un filtre sur \code{latitude IS NULL}, et pas \code{save()}. Comme ça, on écrit seulement les deux colonnes concernées et on n'écrase pas une autre modification de la destination faite au même moment.

\subsection{Countries (GraphQL)}

Le client Countries envoie une requête GraphQL en \code{POST} à \code{countries.trevorblades.com/graphql}. Le code pays est passé dans \code{variables} et n'est jamais collé dans le texte de la requête. C'est la même idée qu'une requête SQL paramétrée, et cela empêche les injections.

\begin{lstlisting}[language=Python,caption={Requête GraphQL paramétrée (\texttt{apps/integrations/services/countries.py}, extrait raccourci).},label={lst:countries}]
COUNTRY_QUERY = """
query Country($code: ID!) {
  country(code: $code) {
    code name native capital emoji currency
    continent { name }
    languages { name native }
  }
}
"""
body = {"query": COUNTRY_QUERY, "variables": {"code": code}}
\end{lstlisting}

En général, un serveur GraphQL répond 200 même quand il y a une erreur. Le client regarde donc la liste \code{errors} avant le champ \code{data}. Si cette liste n'est pas vide, on renvoie une erreur 502. Pour un code inconnu, le service répond \code{{"country": null}}, que nous traduisons en \code{country_not_found} (404).

\subsection{Couche HTTP commune}

Les deux clients utilisent le même module, \code{apps/integrations/services/http.py}. Chaque requête a un délai de 5~secondes et peut être essayée deux fois. Le deuxième essai a lieu seulement si l'erreur est passagère.

\begin{table}[H]
  \centering
  \small
  \caption{Nouvel essai et traduction des échecs.}
  \label{tab:retry}
  \begin{tabularx}{\linewidth}{@{}l c X l@{}}
    \toprule
    \textbf{Situation} & \textbf{Nouvel essai} & \textbf{Exception} & \textbf{HTTP} \\
    \midrule
    Délai dépassé & oui & \code{ExternalServiceTimeout} & 504 \\
    Erreur réseau & oui & \code{ExternalServiceUnavailable} & 503 \\
    Réponse 5xx & oui & \code{ExternalServiceUnavailable} & 503 \\
    Réponse 4xx & non & \code{ExternalServiceUnavailable} & 503 \\
    JSON illisible ou de forme inattendue & non & \code{ExternalServiceBadResponse} & 502 \\
    Ville ou pays inconnu & non & \code{LocationNotFound}, \code{CountryNotFound} & 404 \\
    Séjour hors des 16 jours & non & \code{NoDataAvailable} & 404 \\
    \bottomrule
  \end{tabularx}
\end{table}

On ne réessaie pas après une réponse 4xx, car la même requête serait refusée de la même façon. Les exceptions de \code{apps/integrations/exceptions.py} ne dépendent pas de DRF. Elles contiennent un statut HTTP et un code stable, et chaque API (REST ou GraphQL) les traduit à sa manière. Chaque essai est journalisé avec le service, l'URL sans paramètres, le statut, la durée et le numéro de l'essai.

\subsection{Cache}

Les réponses sont gardées dans le cache de Django, mais seulement si l'analyseur les a acceptées. Une erreur passagère ne reste donc pas bloquée dans le cache jusqu'à l'expiration de l'entrée.

\begin{table}[H]
  \centering
  \small
  \caption{Durées de mise en cache.}
  \label{tab:cache}
  \begin{tabularx}{\linewidth}{@{}l l X@{}}
    \toprule
    \textbf{Donnée} & \textbf{Durée} & \textbf{Clé} \\
    \midrule
    Géocodage & 7 jours & code pays et condensat SHA-256 du nom de ville \\
    Prévisions & 30 minutes & position arrondie à 4 décimales (environ 10~m) \\
    Pays & 24 heures & code pays \\
    \bottomrule
  \end{tabularx}
\end{table}

Nous hachons le nom de la ville pour que la clé soit valide avec n'importe quel moteur de cache, même si l'utilisateur tape des caractères spéciaux. En arrondissant la position, deux destinations très proches partagent la même entrée.

\subsection{Agrégation et dégradation partielle}

L'action \code{GET /api/destinations/{id}/overview/} regroupe dans une seule réponse la destination, son voyage, ses activités, le coût total, la météo et le pays. Les deux services externes sont appelés en parallèle avec un \code{ThreadPoolExecutor}, pendant que le fil principal lit la base.

\begin{lstlisting}[language=Python,caption={Appels en parallèle (\texttt{apps/integrations/services/aggregator.py}, extrait).},label={lst:aggregator}]
with ThreadPoolExecutor(max_workers=len(SOURCES), thread_name_prefix="overview") as pool:
    futures = {
        WEATHER: pool.submit(contextvars.copy_context().run, fetch_forecast, destination),
        COUNTRY: pool.submit(
            contextvars.copy_context().run, countries.get_country, destination.country_code
        ),
    }
    activities = list(destination.activities.all())  # pendant les appels externes
    results, warnings = {}, []
    for source, future in futures.items():
        try:
            results[source] = future.result()
        except IntegrationError as exc:
            results[source] = None
            warnings.append(SourceWarning(source, exc.code, exc.message))
\end{lstlisting}

Ce parallélisme nous a obligés à faire attention à quelques points. D'abord, les fils secondaires ne font aucune requête SQL. Sinon, chaque fil ouvrirait sa propre connexion, qui ne serait jamais fermée, et il ne verrait pas la transaction du fil principal. Le géocodage calcule donc les coordonnées sans les écrire, et c'est le fil principal qui les enregistre après.

Ensuite, chaque tâche tourne dans une copie du contexte (\code{contextvars}). Ses lignes de journal gardent ainsi l'identifiant de la requête en cours.

Enfin, si une source échoue, la réponse ne plante pas et reste un 200. Le bloc de cette source vaut \code{null} et la liste \code{warnings} explique ce qui s'est passé (extrait~\ref{lst:warnings}).

\begin{lstlisting}[language=json,caption={Dégradation partielle quand Open-Meteo est indisponible.},label={lst:warnings}]
"weather": null,
"warnings": [
  {"source": "weather", "code": "external_service_unavailable",
   "message": "Open-Meteo est momentanément indisponible (erreur 503)."}
]
\end{lstlisting}

Cette dégradation partielle sera utile pour le front-end. Même si un fournisseur est en panne, la page d'une destination pourra quand même s'afficher.
